\section{Phương pháp luận (Methodology)}

\subsection{Bộ mã hóa Byte-Level BPE Tối ưu cho Toán học}
Chúng tôi tự cài đặt thuật toán Byte-Level BPE thuần Python (không dùng \texttt{tiktoken} hay \texttt{tokenizers}) khởi tạo từ 256 byte UTF-8 cơ bản và học 16.121 phép ghép (\emph{merges}) để đạt kích thước từ vựng mục tiêu $V = 16{,}384$. Điểm then chốt cho bài toán số học là bộ tiền tách từ (\emph{pre-tokenizer}) sử dụng biểu thức chính quy tách riêng từng chữ số đơn lẻ \texttt{0--9}, đảm bảo mọi số nhiều chữ số (ví dụ \texttt{1234}) luôn được mã hóa nhất quán thành các token đơn lẻ \texttt{['1','2','3','4']} thay vì bị gộp ngẫu nhiên theo tần suất.

Hệ thống cố định 7 token đặc biệt từ ID 0 đến 6: \texttt{<|pad|>}, \texttt{<|bos|>}, \texttt{<|eos|>}, \texttt{<|user|>}, \texttt{<|assistant|>}, \texttt{<think>}, và \texttt{</think>}.

\subsection{Kiến trúc Decoder-Only Transformer}
Mô hình \texttt{DecoderOnlyLM} được xây dựng theo chuẩn Pre-Norm Decoder-Only (họ Llama/Qwen) \cite{touvron2023llama} với cấu hình \emph{Compact} gồm $L_{\text{layers}} = 12$ tầng, $d_{\text{model}} = 512$, $n_{\text{heads}} = 8$ đầu truy vấn, $n_{\text{kv\_heads}} = 4$ đầu Key/Value ($d_{\text{head}} = 64$), $d_{\text{ff}} = 1408$, độ dài ngữ cảnh tối đa $L_{\max} = 2048$, và chia sẻ trọng số Embedding--LM Head (\emph{Weight Tying}), cho tổng số tham số chính xác là $43{,}790{,}848$.

Các thành phần lõi được lập trình tường minh bằng phép toán tensor PyTorch:
\begin{itemize}
    \item \textbf{RMSNorm}: Chuẩn hóa trước mỗi khối con (không dùng bias) giúp ổn định gradient:
    \begin{equation}
        \text{RMSNorm}(\mathbf{x}) = \frac{\mathbf{x}}{\sqrt{\frac{1}{d}\sum_{i=1}^{d} x_i^2 + \epsilon}} \odot \boldsymbol{\gamma}, \quad \epsilon = 10^{-6}.
    \end{equation}
    \item \textbf{RoPE (Rotary Position Embedding)}: Mã hóa vị trí tương đối trực tiếp lên các vector $Q, K$ với cơ số $\theta = 10{,}000$, tương thích tự nhiên với vị trí \texttt{start\_pos} khi giải mã từng bước \cite{su2021roformer}.
    \item \textbf{Grouped-Query Attention (GQA)}: Sử dụng tỷ lệ $2:1$ ($8$ query heads chia sẻ $4$ KV heads), cài đặt thủ công phép tính $\text{softmax}\!\left(\frac{QK^\top}{\sqrt{d_{\text{head}}}} + M_{\text{causal}}\right)V$ (cùng cờ đối chứng \texttt{attn\_backend="sdpa"}).
    \item \textbf{SwiGLU Feed-Forward}: $\text{SwiGLU}(\mathbf{x}) = (\text{SiLU}(\mathbf{x}W_1) \odot (\mathbf{x}W_3))W_2$ với $d_{\text{ff}} = 1408 \approx \frac{8}{3}d_{\text{model}}$.
\end{itemize}

\subsection{Bộ đệm DynamicKVCache và Cơ chế Tái sử dụng Prefix}
Để tránh tính lại phép chiếu Key/Value cho các token cũ khi sinh tự hồi quy, chúng tôi thiết kế \texttt{DynamicKVCache} dạng bộ đệm tensor cấp phát theo từng đoạn (\emph{chunked buffer}) với kích thước khởi tạo \texttt{init\_cap = 128} và tự động mở rộng theo bội số $128$ tối đa đến $2048$. Thiết kế này tránh được chi phí sao chép $O(L^2)$ của \texttt{torch.cat} tại mỗi bước decode đồng thời không lãng phí VRAM cho các chuỗi ngắn. Dung lượng bộ nhớ lý thuyết của KV Cache ở độ chính xác FP16 ($2$ bytes/phần tử) là:
\begin{equation}
    \text{VRAM}_{\text{KV}} = 2 \cdot n_{\text{layers}} \cdot B \cdot n_{\text{kv\_heads}} \cdot L \cdot d_{\text{head}} \cdot 2 \text{ bytes}.
\end{equation}

Đối với Test-Time Scaling ($K$ nhánh song song), hệ thống sử dụng phương thức \texttt{copy\_batch} (\texttt{repeat\_interleave}) để sao chép vật lý trạng thái KV sau pha Prefill thành $K$ nhánh độc lập, đảm bảo cập nhật từng nhánh an toàn tuyệt đối (bên cạnh cấu trúc \texttt{SharedPrefixKVCache} cho tối ưu bộ nhớ mở rộng).

\subsection{Kiểm soát Rò rỉ Dữ liệu và Huấn luyện CoT SFT}
Trước khi tinh chỉnh trên tập \texttt{GSM8K train} (7.473 câu), toàn bộ dữ liệu được đưa qua bộ lọc \texttt{LeakageFilter} 3 tầng đối chiếu với 1.319 câu \texttt{GSM8K test}: (1) khớp chính xác nguyên bản (\emph{Exact}), (2) khớp sau chuẩn hóa chữ thường/dấu câu/số (\emph{Normalized}), và (3) lọc gần trùng lặp với độ tương đồng 5-gram Jaccard $\ge 0.80$.

Dữ liệu CoT SFT được làm sạch thẻ \texttt{<<expr=val>>} và định dạng theo khuôn mẫu chuẩn tách biệt quá trình suy luận và đáp án:
\begin{quote}
\small\texttt{<|bos|><|user|> \{question\} <|assistant|>} \\
\texttt{<think> \{reasoning\_steps\} </think>} \\
\texttt{\#\#\#\# \{final\_answer\}<|eos|>}
\end{quote}
Hàm mất mát Cross-Entropy áp dụng \textbf{Assistant-Only Loss Masking} (gán nhãn \texttt{ignore\_index = -100} cho toàn bộ token thuộc phần câu hỏi của người dùng).

\subsection{Động cơ Test-Time Scaling và Đánh giá GSM8K}
Với mỗi câu hỏi kiểm thử, chúng tôi đánh giá baseline tất định \texttt{Greedy Pass@1} ($T=0$) và sinh một \textbf{shared pool} gồm $K=16$ chuỗi suy luận độc lập ($T=0.7$, $\text{top\_p}=0.95$, $\text{repetition\_penalty}=1.05$) để so sánh công bằng các chiến lược:
\begin{itemize}
    \item \textbf{Majority Voting (\texttt{Maj@4, Maj@8, Maj@16})}: Chuẩn hóa đáp án số trích xuất được về dạng \texttt{Decimal} và chọn đáp án có số phiếu cao nhất (phá hòa bằng log-probability trung bình).
    \item \textbf{Heuristic Best-of-N (\texttt{BoN@8, BoN@16})}: Chấm điểm mỗi chuỗi dựa trên mức độ tuân thủ thẻ \texttt{<think>...</think>} + \texttt{\#\#\#\#}, tỷ lệ phép tính số học hợp lệ (kiểm tra an toàn bằng bộ phân tích biểu thức, không dùng \texttt{eval()}), và log-probability trung bình.
    \item \textbf{Learned ORM Verifier (tùy chọn)}: Mô hình \texttt{OutcomeVerifierModel} gắn đầu tuyến tính lên hidden state tại token đáp án cuối cùng trước \texttt{<|eos|>}, huấn luyện bằng \texttt{BCEWithLogitsLoss} trên các rollout chia theo cấp độ câu hỏi để hỗ trợ \texttt{BoN@K} và \texttt{Weighted Maj@16}.
    \item \textbf{Oracle \texttt{Pass@16}}: Tỷ lệ câu hỏi có ít nhất 1 trong 16 chuỗi đúng (báo cáo như cận trên lý thuyết).
\end{itemize}
